To address climate change, the world must transition to clean energy,
and ocean wave energy is a unique solution to renewably power coastal communities and offshore activities.
Wave energy is currently too expensize to be practical, and system-level optimization can be used to increase its economic viability.
The interdisciplinary nature of wave energy converters requires careful consideration of the optimization formulation and model complexity
to adequately trade off scope, coupling, accuracy, and computation time.
This thesis focuses on the two-body point absorber WEC, inspired by the Reference Model 3 design \cite{RM3}.

% ======================================================================
\subsection{Summary of Project Scope}
\Cref{ch:modeling} considers the system modeling, and develops a novel framework called \texttt{MDOcean} for simulating the hydrodynamics, dynamics, control, structures, and economics of a two-body point absorber WEC.
It uses semi-analytical models to leverage physical and mathematical insight to reduce the computational cost of the simulation,
and validates and benchmarks the code against other models.
Of particular note is the dynamics and controls module, which introduces a method called quasi-linear pseudo-spectral (QLPS) control to efficiently handle constraints and nonlinearities while retaining the effective simplicity of a linear frequency-domain model. 

\Cref{ch:hydrodynamic} elaborates on and extends the hydrodynamics portion of the \texttt{MDOcean} framework, using the matched eigenfunction expansion method to solve the linear potential flow radiation problem for a class of axisymmetric bodies.
It performs numerical experiments to understand the relationship between spatial discretization, series truncation order, nondimensional geometry, accuracy, and runtime.
It further analyzes the block matrix structure for all cases of radially-monotonic geometries,
systematically documents geometry-dependent convergence,
and quantifies accuracy for slanted bodies.

\Cref{ch:optimization} applies the multidisciplinary design optimization technique to formulate, solve, and interpret a WEC optimization.
It develops a set of requirements based on the model and simulation capabilities and design requirements; 
systematically decides the appropriate set of design variables, constraints, parameters, and objectives;
performs single- and multi-objective optimizations and sensitivity analyses;
and interprets the results to yield insight into the design space and the trade-offs between objectives.

\Cref{ch:grid-value}

All chapters directly feed into each other:
the multidisciplinarity, semi-analytical structure, and computational speed of the models in \cref{ch:modeling,ch:hydrodynamic} expose the 
coupling structure, convexity, and smoothness of the design space.
These characteristics in turn inform the optimization formulation decisions of \cref{ch:optimization} and are precisely what enables the optimization to explore the design space so comprehensively and yield improved designs and insight.

\subsection{Key Findings}
Direct comparisons with a boundary element method solver demonstrate that MEEM achieves comparable hydrodynamic coefficients while converging faster with substantially smaller matrices.
This establishes the method's practical computational advantage and strengthens understanding of the link between its mathematical structure and numerical properties.

% ======================================================================
\subsection{Limitations and Future Work}
\subsubsection{Enhancing the Present Scope}

\subsubsection{Vision for Future WEC Design and Optimization}
The same engineering principles and methods used in this thesis can be applied to further improve the design and optimization of WEC systems.
While this work exploits the QCQP structure of the optimal control problem for fixed design,
future work should exploit the fractional-QP structure of the PTO and structural sizing that \cref{sec:qp-convexity} identifies.
This would leave only the hull and structural sizing as a nonconvex design optimization problem, with further structure that can perhaps be exploited even further.
For example, the hull and structural optimizations can be further split up into separate nonconvex subproblems with only linear and quadratic inequality constraints respectively.
Separation reduces the dimension of each problem 
and could potentially mitigate the number of multistart runs required to find a global optimum.
Finally, wrapping the entire optimization in an integer programming framework would allow for the inclusion of discrete design variables, such as the hydrodynamic architecture, number of degrees of freedom, type of PTO, and structural material.
\Cref{fig:nested_future_optim} conceptually demonstrates such an optimization, though a number of practical details remain to be studied.

\begin{figure}
\centering
\label{fig:nested_future_optim}
\includegraphics[width=0.8\textwidth]{figures/nested_future_optim.pdf}
\caption{Conceptual illustration of a nested optimization framework for future WEC design and optimization.}
\end{figure}

The ultimate nesting structure that minimizes the overall solution time may differ from the one proposed in \Cref{fig:nested_future_optim} 
and should be determined based on the relative computational cost of each module and expected number of iterations required for each subproblem (which lowers considerably for convex or analytically solvable subproblems).
These tradeoffs can be systematically explored using MDO, as the optimization architecture remains an  ``asymmetric subspace optimization'' as \citet[Sec. J]{martins_multidisciplinary_2013} describes.
Combination with a multifidelity framework that refines the model complexity as the optimization progresses could enhance the accuracy of the final optimization \cite{wu_gradient-based_2022}.
This could ultimately lead to an optimization framework that chooses the appropriate model fidelity, subproblem structure, and convergence hyperparameters based on dynamically-identified computational cost scalings
that depend on the number of sea states, bodies, controlled degrees of freedom, and hull spatial discretization. 
Such a framework should be created in partnership with WEC developers, energy planners, and coastal communities
to ensure that the process accurately represents real-world design and operation constraints and maximizes the benefits for those ultimately affected.

\subsection{More Meaningful Optimization Metrics}
In this spirit, adjusting the optimization objectives to better align with practical climate-relevant metrics is essential moving forward.
A substantial portion of my PhD research focused on this task, although it is omitted from this dissertation writeup because trustworthy results could not ultimately be obtained in the timeframe available.
A whitepaper will be released describing the work done and the challenges encountered,
and I hope that future researchers will continue to pursue this important task.

% ======================================================================
\subsection{Closing Remarks}